% Ch15 WS3 -- titration curves, indicators, and a 10-point FRQ. Blocks 4-5.
\documentclass{apchem-worksheet}

\apcunitword{Chapter}
\apcunit{15}
\apcunittitle{Acid--Base Equilibria}
\apcdoctitle{Worksheet 3 \textbullet\ Titration Curves, Indicators \& FRQ}
\apcsubtitle{Zumdahl \S15.4--15.5 \textbullet\ name the major species before choosing a method}

\begin{document}
\wsheader[$K_a$: \ce{HC2H3O2} $1.8\times10^{-5}$ (p$K_a$ 4.74) \textbullet\
\ce{HOCl} $3.5\times10^{-8}$ (p$K_a$ 7.46). $K_b(\ce{NH3}) =
1.8\times10^{-5}$, $K_a(\ce{NH4+}) = 5.6\times10^{-10}$ (p$K_a$ 9.25).
$K_w = 1.0\times10^{-14}$.]

\problem[]{For a weak acid titrated with a strong base, name the major
species and the method used in each region:}
\begin{parts}
  \item Before any base is added:
        \answerblank[7cm]{\ce{HA} only --- weak-acid ICE}
  \item Between the start and equivalence:
        \answerblank[7cm]{\ce{HA} and \ce{A-} --- buffer, use H--H}
  \item At equivalence:
        \answerblank[7cm]{\ce{A-} only --- weak-base ICE with $K_b$}
  \item After equivalence:
        \answerblank[7cm]{excess \ce{OH-} --- strong base alone}
\end{parts}

\problem[]{\SI{25.0}{\milli\litre} of \SI{0.100}{\Molar} \ce{HCl} is
titrated with \SI{0.100}{\Molar} \ce{NaOH}.}
\begin{parts}
  \item Volume of \ce{NaOH} at equivalence:
        \answerblank[3cm]{\SI{25.0}{\milli\litre}}
  \item pH before any base is added: \answerblank[2.2cm]{1.00}
  \item pH after \SI{10.0}{\milli\litre}: \workspace[2cm]
        \keyonly{\begin{apcanswer}\ce{H+} left $= 2.50 - 1.00 =
        \SI{1.50}{\milli\mole}$ in \SI{35.0}{\milli\litre};
        $[\ce{H+}] = 0.0429$; pH $= \mathbf{1.37}$\end{apcanswer}}
  \item pH at equivalence, with a reason: \answerlines[2]{7.00 --- the only
        species left are \ce{Na+}, \ce{Cl-} and water, and neither ion
        reacts with water.}
  \item pH after \SI{35.0}{\milli\litre}: \workspace[2cm]
        \keyonly{\begin{apcanswer}excess \ce{OH-} $= 3.50 - 2.50 =
        \SI{1.00}{\milli\mole}$ in \SI{60.0}{\milli\litre};
        $[\ce{OH-}] = 0.0167$; pOH $= 1.78$;
        pH $= \mathbf{12.22}$\end{apcanswer}}
\end{parts}

\problem[]{\SI{25.00}{\milli\litre} of \SI{0.100}{\Molar} \ce{HC2H3O2} is
titrated with \SI{0.100}{\Molar} \ce{NaOH}. Complete the table.}
\begin{parts}
  \item \SI{0.00}{\milli\litre} (weak acid alone):
        \answerblank[2.4cm]{2.87}
  \item \SI{5.00}{\milli\litre}: \workspace[2cm]
        \keyonly{\begin{apcanswer}\ce{HA} $= 2.500-0.500 = 2.000$,
        \ce{A-} $= 0.500$ mmol;
        pH $= 4.74 + \log(0.500/2.000) = 4.74 - 0.60 =
        \mathbf{4.14}$\end{apcanswer}}
  \item \SI{12.50}{\milli\litre}, and say what is special about it:
        \answerlines[2]{The halfway point --- $[\ce{HA}] = [\ce{A^-}]$, so
        pH $=$ p$K_a$ $= 4.74$.}
  \item \SI{20.00}{\milli\litre}: \workspace[2cm]
        \keyonly{\begin{apcanswer}\ce{HA} $= 0.500$, \ce{A-} $= 2.000$ mmol;
        pH $= 4.74 + \log(2.000/0.500) = 4.74 + 0.60 =
        \mathbf{5.34}$\end{apcanswer}}
  \item \SI{25.00}{\milli\litre} (equivalence): \workspace[2.6cm]
        \keyonly{\begin{apcanswer}$[\ce{C2H3O2^-}] = 2.500/50.00 = 0.0500$;
        $K_b = 5.6\times10^{-10}$;
        $x = 5.3\times10^{-6}$; pOH $= 5.28$;
        pH $= \mathbf{8.72}$\end{apcanswer}}
  \item \SI{30.00}{\milli\litre}: \workspace[2cm]
        \keyonly{\begin{apcanswer}excess \ce{OH-} $= 0.500$ mmol in
        \SI{55.00}{\milli\litre}; $[\ce{OH-}] = 9.09\times10^{-3}$;
        pOH $= 2.04$; pH $= \mathbf{11.96}$\end{apcanswer}}
\end{parts}

\problem[]{Compare the two titrations in problems 2 and 3. Both use the same
volumes and concentrations.}
\begin{parts}
  \item Are the equivalence \emph{volumes} the same? Explain.
        \answerlines[2]{Yes --- \SI{25.0}{\milli\litre} in both. Equivalence
        volume depends only on stoichiometry and concentration, not on acid
        strength.}
  \item Are the equivalence \emph{pH values} the same? Explain.
        \answerlines[3]{No: 7.00 for \ce{HCl}, 8.72 for acetic acid. The
        acetic acid titration leaves \ce{C2H3O2^-}, a weak base that
        hydrolyses to give \ce{OH-}; the \ce{HCl} titration leaves only
        inert \ce{Cl-}.}
  \item Which curve has a flat buffer region, and why?
        \answerlines[2]{The acetic acid curve. Between the start and
        equivalence it contains both \ce{HC2H3O2} and \ce{C2H3O2^-}, so
        added base only changes their ratio, and the pH depends on the log
        of that ratio.}
\end{parts}

\problem[]{\SI{100.0}{\milli\litre} of \SI{0.050}{\Molar} \ce{NH3} is
titrated with \SI{0.10}{\Molar} \ce{HCl}.}
\begin{parts}
  \item Volume at equivalence: \answerblank[3cm]{\SI{50.0}{\milli\litre}}
  \item pH at the halfway point, with a one-line reason:
        \answerlines[2]{\SI{25.0}{\milli\litre}: $[\ce{NH3}] =
        [\ce{NH4+}]$, so pH $=$ p$K_a(\ce{NH4+}) = 9.25$.}
  \item pH at equivalence: \workspace[2.6cm]
        \keyonly{\begin{apcanswer}$[\ce{NH4+}] = 5.0/150.0 = 0.0333$;
        $K_a = 5.6\times10^{-10}$;
        $x = \sqrt{(5.6\times10^{-10})(0.0333)} = 4.3\times10^{-6}$;
        pH $= \mathbf{5.36}$\end{apcanswer}}
  \item Why is this equivalence point acidic?
        \answerlines[2]{All the \ce{NH3} has been converted to \ce{NH4+},
        the conjugate acid of a weak base, which donates a proton to water.}
\end{parts}

\problem[]{Indicators.}
\begin{parts}
  \item An indicator is itself a \answerblank[3.4cm]{weak acid}, whose
        \ce{HIn} and \ce{In-} forms differ in
        \answerblank[2cm]{color}.
  \item An indicator with $K_a = 1.0\times10^{-5}$ changes color over
        what pH range? \answerblank[3cm]{about 4 to 6}
  \item Distinguish the end point from the equivalence point.
        \answerlines[2]{The equivalence point is fixed by reaction
        stoichiometry; the end point is where the indicator visibly changes
        color. A well-chosen indicator makes them nearly coincide.}
  \item Which indicator suits the titration in problem 3
        (equivalence 8.72): methyl red (4.4--6.2) or phenolphthalein
        (8.2--10.0)? \answerblank[4cm]{phenolphthalein}
  \item Which suits problem 5 (equivalence 5.36)?
        \answerblank[4cm]{methyl red}
\end{parts}

\problem[]{\textbf{FRQ (10 points).} A \SI{20.00}{\milli\litre} sample of a
solution of an unknown monoprotic weak acid \ce{HA} is titrated with
\SI{0.1000}{\Molar} \ce{NaOH}. The equivalence point is reached after
\SI{22.40}{\milli\litre}. The pH after \SI{11.20}{\milli\litre} of base is
4.85.}
\begin{parts}
  \item Calculate the concentration of the acid. \workspace[2.4cm]
        \keyonly{\begin{apcanswer}$n = (22.40)(0.1000) =
        \SI{2.240}{\milli\mole}$;
        $C = 2.240/20.00 = \mathbf{0.1120}~\mathrm{M}$\end{apcanswer}}
  \item Explain why the pH at \SI{11.20}{\milli\litre} gives p$K_a$
        directly. \answerlines[3]{That volume is exactly half the
        equivalence volume, so exactly half the \ce{HA} has been converted
        to \ce{A-} and $[\ce{HA}] = [\ce{A^-}]$. The log term in
        Henderson--Hasselbalch is then zero and pH $=$ p$K_a$.}
  \item Determine $K_a$. \answerblank[3.4cm]{$1.4\times10^{-5}$}
  \item Calculate the pH at the equivalence point. \workspace[3cm]
        \keyonly{\begin{apcanswer}$[\ce{A^-}] = 2.240/42.40 = 0.05283$;
        $K_b = 1.0\times10^{-14}/1.4\times10^{-5} = 7.1\times10^{-10}$;
        $x = \sqrt{(7.1\times10^{-10})(0.05283)} = 6.1\times10^{-6}$;
        pOH $= 5.21$; pH $= \mathbf{8.79}$\end{apcanswer}}
  \item Choose an indicator for this titration from methyl red (4.4--6.2),
        bromthymol blue (6.0--7.6) and phenolphthalein (8.2--10.0), and
        justify. \answerlines[2]{Phenolphthalein --- its range brackets the
        equivalence pH of 8.79, so its end point falls essentially at
        equivalence. The other two would change color while acid remained.}
\end{parts}

\begin{apcnote}[Rubric --- score yourself, 10 points]
\textbf{(a) 2 pts} 1 mmol of base, 1 concentration \textbullet\
\textbf{(b) 2 pts} 1 identifies it as half the equivalence volume, 1 states
$[\ce{HA}] = [\ce{A^-}]$ makes the log term zero --- asserting
``pH $=$ p$K_a$ at halfway'' with no reasoning earns 1 \textbullet\
\textbf{(c) 1 pt} $K_a = 10^{-4.85} = 1.4\times10^{-5}$ \textbullet\
\textbf{(d) 3 pts} 1 dilution to the total \SI{42.40}{\milli\litre},
1 $K_b$ from $K_w/K_a$, 1 pH via pOH --- using $[\ce{A^-}] = 0.1120$
(forgetting dilution) caps this part at 2 \textbullet\
\textbf{(e) 2 pts} 1 phenolphthalein, 1 justification referring to the
equivalence pH.
\end{apcnote}

\begin{spiralstrip}{Chapter 15 \textbullet\ blocks 1--3}
  \item Step 1 of the buffer method:
        \answerblank[5.7cm]{stoichiometry to completion}
  \item Optimal buffering ratio: \answerblank[1.6cm]{1}
  \item Useful buffer range: \answerblank[2.4cm]{p$K_a \pm 1$}
  \item pH of \SI{0.20}{\Molar}/\SI{0.20}{\Molar} acetate buffer:
        \answerblank[1.8cm]{4.74}
  \item Diluting a buffer changes capacity how?
        \answerblank[2.4cm]{lowers it}
\end{spiralstrip}

\end{document}
